\section{تحديد الثوابت الميكانيكية ومصفوفات القصور الذاتي للوصلات}
\label{sec:mass_inertia_parameters}

يتطلب حساب النموذج الديناميكي معرفة دقيقة بالثوابت الفيزيائية والميكانيكية لكل وصلة من وصلات الروبوت. يتضمن ذلك أطوال الوصلات، وكتلها، وتحديد متجهات مراكز الثقل، بالإضافة إلى حساب مصفوفات القصور الذاتي المتكاملة (\LR{Pseudo-Inertia Matrices}) التي تُشكل النواة الرياضية في صياغة لاغرانج-أويلر.

\subsection{المعايير الهندسية والكتلية للروبوت}
بناءً على التصميم الميكانيكي الفعلي لروبوت (\LR{SCARA}) المدروس، تم تحديد قيم الأطوال والكتل الثابتة للوصلات الثلاث، بالإضافة إلى دمج كتلة الحمل المتغير $(m_L)$ عند نهاية الوصلة الثالثة كما يلي:
\begin{itemize}
    \item أطوال الوصلات: $L_1 = 0.45 \text{ m}$، و $L_2 = 0.35 \text{ m}$، و $L_3 = 0.1 \text{ m}$.
    \item كتل الوصلات الأصلية (بدون حمل): $m_1 = 14 \text{ kg}$، و $m_2 = 18 \text{ kg}$، و $m_3 = 2 \text{ kg}$.
    \item الكتلة المركبة الإجمالية للوصلة الثالثة مع الحمل: $m_3^* = m_3 + m_L$.
\end{itemize}

\subsection{متجهات مراكز الثقل (\LR{Center of Mass Vectors})}
بفرض أن الوصلتين الأولى والثانية عبارة عن قضبان متجانسة الحجم والكثافة، فإن مركز ثقلهما الهندسي يقع في المنتصف تماماً على امتداد المحور الطولي الفعال. أما الوصلة الثالثة، فنظراً لكونها الجزء المسؤول عن حمل الأوزان المتغيرة، فإن مركز ثقلها الكلي ينزاح لحظياً بناءً على قيمة الحمل المرفوع $(m_L)$. 

باستخدام قانون العزوم السكونية لمركز الكتل المشترك :
\begin{equation}
R_c = \frac{\sum m_i r_i}{\sum m_i}
\end{equation}

تم تحديد متجهات مراكز الثقل $r_{p1}$ و $r_{p2}$ و $r_{p3}^*$ لكل وصلة منسوبة إلى جملة إحداثياتها المحلية كما يلي:

\begin{equation}
r_{p1} = \begin{bmatrix} \frac{L_1}{2} \\ 0 \\ 0 \end{bmatrix} , \quad 
r_{p2} = \begin{bmatrix} \frac{L_2}{2} \\ 0 \\ 0 \end{bmatrix} , \quad 
r_{p3}^* = \begin{bmatrix} 0 \\ 0 \\ z_{c3}^* \end{bmatrix}
\label{eq:com_vectors}
\end{equation}

حيث يمثل الإحداثي $z_{c3}^*$ موضع مركز الثقل المشترك للوصلة الثالثة مع الحمل على المحور العمودي، ويُحسب لحظياً عبر العلاقة التالية:
\begin{equation}
z_{c3}^* = \frac{m_3 \left(\frac{L_3}{2}\right) + m_L L_3}{m_3 + m_L}
\label{eq:zc3_star}
\end{equation}

\subsection{مصفوفات القصور الذاتي المتكاملة للوصلات (\LR{Pseudo-Inertia Tensors})}
في صياغة لاغرانج-أويلر، يتم دمج مصفوفة عزم القصور الذاتي الدوراني القياسية $I_i$ مع متجهات مركز الثقل والكتلة ضمن مصفوفة واحدة ممتدة بأبعاد $4 \times 4$ تُعرف بالمصفوفة المتكاملة $J_i$. تُعرّف هذه المصفوفة عموماً بالصيغة التالية:

\begin{equation}
J_i = \begin{bmatrix} 
\int x^2 dm & \int xy dm & \int xz dm & \int x dm \\ 
\int xy dm & \int y^2 dm & \int yz dm & \int y dm \\ 
\int xz dm & \int yz dm & \int z^2 dm & \int z dm \\ 
\int x dm & \int y dm & \int z dm & \int dm 
\end{bmatrix}
\end{equation}

بتطبيق هذا المفهوم وحساب التكاملات الحجمية للوصلات بناءً على توزيع الكتل المحدد ومراكز الثقل في المعادلات (\ref{eq:com_vectors})، تم استنتاج مصفوفات القصور الذاتي المتكاملة الجاهزة للتعويض المباشر في معادلات الحركة:

\textbf{1. مصفوفة القصور الذاتي للوصلة الأولى ($J_1$):}
\begin{equation}
J_1 = \begin{bmatrix}
\frac{1}{3}m_1 L_1^2 & 0 & 0 & \frac{L_1}{2}m_1 \\
0 & 0 & 0 & 0 \\
0 & 0 & 0 & 0 \\
\frac{L_1}{2}m_1 & 0 & 0 & m_1
\end{bmatrix}
\end{equation}

\textbf{2. مصفوفة القصور الذاتي للوصلة الثانية ($J_2$):}
\begin{equation}
J_2 = \begin{bmatrix}
\frac{1}{3}m_2 L_2^2 & 0 & 0 & \frac{L_2}{2}m_2 \\
0 & 0 & 0 & 0 \\
0 & 0 & 0 & 0 \\
\frac{L_2}{2}m_2 & 0 & 0 & m_2
\end{bmatrix}
\end{equation}

\textbf{3. مصفوفة القصور الذاتي المركبة للوصلة الثالثة مع الحمل ($J_3$):}
نظراً لتغير مركز الثقل، تم استخدام نظرية المحاور المتوازية لشتاينر (\LR{Steiner's Theorem}) لحساب انزياح عزم القصور الذاتي للاسطوانة والحمل عن المركز الجديد، حيث تعطى بالصيغة العامة $I_{new} = I_{center} + md^2$. بناءً على ذلك، تم صياغة المصفوفة المتكاملة المحدثة للحمل والوصلة الثالثة ككتلة واحدة بالشكل التالي:

\begin{equation}
J_3 = \begin{bmatrix}
0 & 0 & 0 & 0 \\
0 & 0 & 0 & 0 \\
0 & 0 & J_{3z} & -\left(m_3 \frac{L_3}{2} + m_L L_3\right) \\
0 & 0 & -\left(m_3 \frac{L_3}{2} + m_L L_3\right) & m_3 + m_L
\end{bmatrix}
\end{equation}

حيث يمثل العنصر $J_{3z}$ عزم القصور الذاتي الإجمالي الفعال والمصحح حول المحور العمودي، ويُحسب بالاعتماد على انزياح الاسطوانة وانزياح الحمل عن المركز المشترك :
\begin{equation}
J_{3z} = \frac{1}{3}m_3 L_3^2 + 2m_3 h_3^2 + 2m_L h_L^2 - \frac{1}{2}m_3 r^2
\end{equation}

حيث أن $h_3 = \left| \frac{L_3}{2} - z_{c3}^* \right|$ يمثل انزياح مركز ثقل الوصلة الأصلية، و $h_L = \left| L_3 - z_{c3}^* \right|$ يمثل انزياح مركز ثقل الحمل، و $r$ هو نصف قطر اسطوانة الوصلة الثالثة.